\talkslide{limitations}{Limitations and future work}{Timing: provisional}{The published paper gives an equilibrium argument for an idealized tabular formulation, but that is distinct from proving convergence of neural training or certifying all possible attacks.}{
\begin{tikzpicture}[remember picture,overlay,
 every node/.style={text=Ink,align=left,inner sep=0pt,font=\normalsize}]
\begin{scope}[shift={(current page.south west)}]
 \node[anchor=north west,font=\normalsize\bfseries] at (.65,7.05) {Limits of the defense};
 \node[anchor=north west,text width=14.2cm] at (1.1,6.6)
  {No robustness certificate or known ground-truth equilibrium\par
   Finding one failure is easier than defending against \textbf{every attack}};
 \draw[Muted!35,line width=.5pt] (.65,5.55)--(15.1,5.55);
 \node[anchor=north west,font=\normalsize\bfseries] at (.65,5.2) {Broader attack coverage};
 \node[anchor=north west,text width=14.2cm] at (1.1,4.75)
  {Augment learned attacks with human-defined collections and strategies\par
   Known attack families are not manually integrated into the current game};
 \draw[Muted!35,line width=.5pt] (.65,3.7)--(15.1,3.7);
 \node[anchor=north west,font=\normalsize\bfseries] at (.65,3.35) {Prompt injection on agents};
 \node[anchor=north west,text width=14.2cm] at (1.1,2.9)
  {Attacks through \textbf{untrusted content} in an agent's environment};
 \draw[Muted!35,line width=.5pt] (.65,2.2)--(15.1,2.2);
 \node[anchor=north west,font=\normalsize\bfseries] at (.65,1.85) {Other modalities};
 \node[anchor=north west,text width=14.2cm] at (1.1,1.4)
  {Extensions to multimodal and non-textual inputs};
\end{scope}
\end{tikzpicture}

}
\speak{The published paper gives an equilibrium argument for an idealized tabular formulation, but that is distinct from proving convergence of neural training or certifying all possible attacks. The attacker only covers part of the space, and a judge can misread faithfulness or safety. Known attack families are not exhaustively injected into the training game. An attacker only needs one failure, while a defender must preserve broad useful behavior. Finally, the recipe costs more than a single-model defense, although it produces a reusable attacker as well. Source detail: AdvGame, Conclusion and Appendices B--C. One longer-term hope is that the learned attacker discovers and subsumes earlier attack strategies, adapting to the current defender. That is an aspiration, not something the experiments establish. A practical future direction is to augment learned generation with human-defined attack collections and strategies. We did not exhaustively integrate these families in the evaluated recipe. In an agent setting, an attacker can modify untrusted content the model encounters while pursuing a legitimate task. The defender needs to complete that task without following the injected instruction. The game would need constraints on what the attacker can edit and an evaluation of both task completion and unsafe actions. A second direction is extending attacks and defenses to multimodal and non-textual inputs. These are future research directions, not experiments established by this paper.}
